\section{Experimental design}\label{sec:setup}

\subsection{Streams and protocol}
Each realisation (seed) generates its own channel trajectories, calibration data, training data and test streams; seeds are the unit of replication. A test stream covers 24 hours (8\,640 windows). Benign streams are run at drift scales $d\in\{0,0.25,1,2,4\}$; $d=0.25$ represents a well-stabilised link with misalignment growth of about 0.25\%/h. For each scale and each attack of Table~\ref{tab:attacks}, a further stream carries the attack for two hours (720 windows) starting at a uniformly random window between hour 4 and hour 20. The same channel trajectory is also sampled without the attack, which gives a paired honest copy of the attacked windows. Counts are drawn per photon-number component ($n=0$, $n=1$, $n\geq2$) and then thinned by Bob's basis choice, so every block records its true single-photon detections and the number of sifted key bits Eve knows. The finite-key analysis runs on every complete block of 60 windows, on a fixed grid of the stream, that lies inside the attack, and on all 144 blocks of each benign stream.

\subsection{Detectors}
Thirteen detectors share one interface: a score per window and a threshold.
\begin{itemize}
\item \textbf{QBER}: the sifted $Z$-basis QBER of the signal intensity against the fixed 11\% limit used as the baseline in \cite{isha2026}.
\item \textbf{Band}: a two-sided physics band. The signal QBER and the decoy-to-signal gain ratio per unit monitored intensity are standardised with the median and MAD of a drift-free benign stream, and the score is the larger absolute $z$-value. It alarms on QBER that is too low as well as too high.
\item \textbf{DCS}, \textbf{CUSUM} and \textbf{Block} (DCS-block): Section~\ref{sec:dcs}. CUSUM and Block are sequential: their window scores are not independent, so their per-window FPR is not directly comparable with that of the single-window detectors, and alarm episodes and time to first false alarm are the more relevant measures for them. The CUSUM does not reset after an alarm, so it stays in alarm while the evidence persists; its per-window TPR measures how much of the attack it covers, not how fast it detects. A CUSUM on the Band features would inherit Band's drift failure (Section~\ref{sec:fpr-results}) and is not included. \textbf{DCS-X} is the variant that uses only test-basis errors (Section~\ref{sec:dcs}).
\item \textbf{Band-roll}: the Band features with an adaptive reference, each minus its causal rolling median over the previous hour (360 windows), standardised and thresholded on drifting benign data. It represents the rolling baselines operators use; a shorter reference absorbs a two-hour attack after onset.
\item \textbf{IF-stat}, \textbf{IF-drift} and \textbf{IF-res}: isolation forests \cite{liu2008} with 200 trees. IF-stat and IF-drift use nine features: the logarithm of the gain per intensity and the sifted QBER per basis and intensity. IF-stat is fitted on one drift-free benign stream, IF-drift on three benign streams with $d\in\{0.5,1,2\}$. IF-res is fitted on the four standardised DCS residuals of the drift-free stream.
\item \textbf{RF-stat}, \textbf{RF-DR} and \textbf{RF-res}: random forests \cite{breiman2001} (200 trees, minimum leaf size 2) trained as binary attack-versus-benign classifiers in the style of the literature \cite{isha2026,alkuwari2026}. The attack class holds six streams of 2\,160 windows (IR with $f\in\{0.05,0.1,0.25,0.5\}$, PNS, Blind-naive). Added loss is a benign event from the protocol's point of view and is not used as an attack label. RF-stat sees one drift-free benign stream and no drift. RF-DR sees six benign streams and draws a drift scale from $U(0,4)$ for every training stream, the domain randomisation of \cite{mohamed2026}. RF-res uses the same data as RF-DR with the standardised DCS residuals as features. IR at $f=0.05$, PNS and Blind-naive are therefore in the training distribution of the supervised detectors (marked with $^*$ in the tables); no detector sees the other blinding variants.
\end{itemize}

\subsection{Thresholds}
Every threshold except the fixed QBER limit is the $(1-\alpha)$ quantile, $\alpha=10^{-3}$, of the detector's scores on three calibration streams (25\,920 windows), with randomised tie-breaking: a score equal to the threshold alarms with the probability that makes the calibration alarm rate exactly $\alpha$. Without it, a random forest whose calibration scores pile up at 0 or 1 would either never or always alarm. Two calibration sets exist. The drift-free set is what an operator has after commissioning; it is the default for every detector not built to use drift. IF-drift, RF-DR, RF-res and Band-roll are drift-aware by design and default to the drifting set ($d\in\{0.5,1,2\}$, marked $^\dagger$). Each detector is also evaluated with the other set as a control. Test drift $d=4$ lies outside both calibration sets; for RF-DR and RF-res it lies at the edge of the training range.

\emph{Seed history.} The benchmark was developed on seeds 0--9, which were run several times while the design changed in response to review: four instead of seven residuals, the Block, DCS-X and Band-roll detectors, the basis-dependent drift terms, the monitor bias, the attackers Blind-lag, Blind-sub and Blind-partial, and the block-wise release policy. The design was then frozen, and every result reported here comes from seeds not used during development: seeds 10--49 for the main benchmark, 500--509 for DCS+F, 600--609 for the stress test and 700--709 for the sensitivity runs. DCS+P, the first floor variant, failed during development and is reported from its development seeds 200--209. The Block detector and the floor variants remain post-hoc in the sense that they were motivated by development results. IR at $f=0.2$ lies between two training strengths of the supervised detectors and is marked as interpolated ($^\ddagger$).

\subsection{Metrics and statistics}
For each detector the benchmark records:
\begin{itemize}
\item the per-window false-alarm rate on benign streams (FPR, expected alarm probability under the randomised threshold) and the number of alarm episodes per day, an episode being a maximal run of consecutive alarmed windows;
\item the true-positive rate on attacked windows (TPR);
\item the paired ROC AUC of the 720 attacked windows against the same 720 windows of the honest copy, which removes time-of-day confounding between attacked and non-attacked parts of a stream. It is the Mann--Whitney statistic of two time-matched samples and does not exploit the pairing further;
\item the delay from attack onset to the first alarm, in windows;
\item the TPR at a matched FPR, with the threshold set at the $(1-\alpha)$ quantile of the detector's own benign scores at the same drift scale. This oracle threshold is estimated in-sample from about nine exceedances, is not available to an operator, and is one-sided, so it cannot register attacks that make windows look more normal; it separates a detector's discriminating power from its calibration;
\item the time to the first false alarm on benign streams, an empirical in-control run length, which is the more relevant measure for the sequential detectors.
\end{itemize}
An AUC below 0.5 means the attack makes windows look \emph{more} normal than honest operation. A one-sided detector does not alarm on such windows, but the attack is still separable, and we report it as such. For each block the benchmark records $\ell$, $n_Z$, $K_E$, the breach indicator (\ref{eq:breach}), the bound checks $s_{Z,1}\leq$ true single-photon count in $Z$ and $s_{X,1}\leq$ its $X$ counterpart, and the phase-error check $v_{X,1}\geq$ true single-photon errors in $X$. It also records the same quantities with nominal instead of monitored intensities, and runs the same analysis on the paired honest copy. To show what an alarm buys, we apply a block-wise release policy: a block is withheld by a detector if that detector alarmed in any of its windows, and release resumes with the next block. Applied to the paired honest copy, the same policy gives the cost in withheld honest blocks, so that withholding by detection can be told apart from withholding by false alarms.

Means over the \nSeeds{} seeds carry 95\% BCa bootstrap intervals over seeds \cite{efron1979} ($10^4$ resamples, a fixed generator per interval). When all seeds give the same value the interval is reported as degenerate, and when the BCa correction is undefined the percentile interval is used instead. Because a signed AUC averaged over seeds can hide separations of opposite sign, we also report the mean of $|\mathrm{AUC}-0.5|$ per seed for the attacks claimed to be invisible. Differences between the DCS and another detector are paired by seed and bootstrapped the same way. We report effect sizes with intervals rather than significance tests. Since one seed corresponds to one link, we also give the largest per-seed false-alarm rate and the share of seeds above $2\alpha$. One seed of the full benchmark takes about \runtime{} s on four CPU cores.
